% Pista ana-24s-q5 · Anàlisi
% Procedència: PAU Catalunya 2024 · Convocatòria de setembre · Sèrie 3 · Qüestió 5
\documentclass[11pt,a4paper]{article}
\usepackage{pista}
\pistainfo{Catalunya 2024}{Convocatòria de setembre}{Sèrie 3}

\begin{document}

\section*{\textcolor{blauPista}{Qüestió 5}}

\noindent Per a cada punt $(x, y)$ de la corba $y = e^{-2x}$, amb $x > 0$ i $y > 0$, considereu el rectangle amb vèrtexs als punts $(0, 0)$, $(x, 0)$, $(0, y)$ i $(x, y)$.

\subsection*{\textit{a)} \normalfont Comproveu que, d'entre tots aquests rectangles, el que té $x = \dfrac{1}{2}$ és el d'àrea màxima. Quin és el valor d'aquesta àrea? \hfill\textnormal{[1,5 punts]}}

\pista{El rectangle té base $x$ i alçada $y = e^{-2x}$. Per tant, l'àrea com a funció d'$x$ és $A(x) = x \cdot e^{-2x}$. És un problema d'optimització: deriva $A(x)$ (recorda que estàs derivant un producte de funcions); després, resol l'equació $A'(x) = 0$, i, finalment, comprova que $x = \dfrac{1}{2}$ és efectivament un màxim. Finalment, avalua $A\!\left(\dfrac{1}{2}\right)$ per a obtenir el valor de l'àrea màxima. Aquest és un tipus de problemes que JA ET DIUEN LA SOLUCIÓ, per tant, és més senzill que si et demanen: "troba el rectangle d'àrea màxima".}

\subsection*{\textit{b)} \normalfont Calculeu l'equació de la recta tangent a la funció $y = e^{-2x}$ en el punt d'abscissa $x = 0$, i el seu punt de tall amb l'eix de les abscisses. \hfill\textnormal{[1 punt]}}

\pista{Aplica la fórmula de la recta tangent, $y - f(0) = f'(0)(x - 0)$. Avalua $f(0)$ i $f'(0)$ (la funció derivada, ja la tens de l'apartat anterior). Un cop tinguis l'equació de la recta tangent, per trobar el punt de tall amb l'eix $OX$ imposa $y = 0$ i resol per a $x$. Recorda que la teva resposta final ha de ser un punt amb les dues coordenades.}

\end{document}
